Deleting a lower-trust source does not necessarily remove information copied into an agent's derived memory. Provenance-aware deletion can follow recorded links, but missing links change both deletion reach and a defender's evidence. MemoryFence separates true ancestry for evaluation from operational recorded ancestry. Its preserved bounded model study contains three attack chains with two descendants each: recorded cascade leaves 0/6, 4/6 and 6/6 active attack descendants at nominal edge-retention probabilities 1.0, 0.75 and 0.5. Eighteen complete-lineage attack episodes have zero exact attacker-objective tool effects; this does not establish protection under missing lineage. We extend the framework with opt-in observable- lineage containment and evaluate seven strategies in 560 paired deterministic runs over 80 configurations from 16 graph families. In four held-out complete-lineage attack cases, selective re-derivation with or without containment preserves 8/8 unrelated derived facts versus cascade's 2/8, while exposing no poisoned context. Marked partial lineage allows containment but loses all 8/8 benign facts; silent loss exposes poisoned context in 3/4 held-out cases. These results characterize software-level lifecycle and utility tradeoffs, not new LLM attack success rates.
